% DP86: complete LO collinear kernels with explicit spin/color contractions.
\begin{derivation}[从\eqnrefs{eq:qcd-collinear-virtuality-dp68}到\eqnrefs{eq:qcd-lo-splitting-kernels-rest-dp83}]
把共线极限中出现的横向四矢量写成 $k_\perp^\mu$，并定义正的欧几里得横向大小
\begin{equation*}
 \kappa_\perp^2\equiv-k_\perp^2>0.
\end{equation*}
取两条类光参考方向 $P^\mu,n^\mu$，满足 $P^2=n^2=0$、$P\cdot n\neq0$。两个无质量子动量统一参数化为
\begin{align*}
 p^\mu
 &=zP^\mu+k_\perp^\mu
 +\frac{\kappa_\perp^2}{2zP\cdot n}\,n^\mu,\\
 k^\mu
 &=(1-z)P^\mu-k_\perp^\mu
 +\frac{\kappa_\perp^2}{2(1-z)P\cdot n}\,n^\mu,
\end{align*}
其中 $P\cdot k_\perp=n\cdot k_\perp=0$。直接平方可得
\begin{align*}
 p^2&=0,\qquad k^2=0,\\
 2p\cdot k
 &=\frac{\kappa_\perp^2}{z(1-z)},\\
 s_{pk}\equiv(p+k)^2
 &=\frac{\kappa_\perp^2}{z(1-z)},
\end{align*}
这与正文\eqnrefs{eq:qcd-collinear-virtuality-dp68} 相同。近在壳母线的传播子奇异因子 $1/s_{pk}$ 可以与子分裂顶角分离，而硬过程只解析到类光母方向 $P$。

先算 $g(P)\to q(p)\bar q(k)$。母胶子的物理偏振投影取轴向规范形式
\begin{equation*}
 d_{\mu\nu}(P,n)
 =-\eta_{\mu\nu}
 +\frac{P_\mu n_\nu+n_\mu P_\nu}{P\cdot n},
 \qquad
 P^\mu d_{\mu\nu}=n^\mu d_{\mu\nu}=0.
\end{equation*}
四维中母胶子有两个物理偏振，因此未极化分裂要乘 $1/2$。颜色平均给
\begin{align*}
 \frac1{N_c^2-1}\sum_{a,i,j}
 (T^a)_{ij}(T^a)_{ji}
 &=\frac1{N_c^2-1}\sum_a\Tr(T^aT^a)\\
 &=T_F.
\end{align*}
末态自旋求和后，分裂顶角的 Lorentz 分子为
\begin{align*}
 \mathcal N_{q\bar q}
 &\equiv\frac{T_F}{2}
 d_{\mu\nu}(P,n)
 \Tr\!\left(\slashed p\gamma^\mu\slashed k\gamma^\nu\right)\\
 &=2T_F\,d_{\mu\nu}(P,n)
 \left[p^\mu k^\nu+p^\nu k^\mu
 -\eta^{\mu\nu}p\cdot k\right],
\end{align*}
其中使用
\begin{equation*}
 \Tr(\gamma^\alpha\gamma^\mu\gamma^\beta\gamma^\nu)
 =4\left(
 \eta^{\alpha\mu}\eta^{\beta\nu}
 -\eta^{\alpha\beta}\eta^{\mu\nu}
 +\eta^{\alpha\nu}\eta^{\mu\beta}
 \right).
\end{equation*}
先收缩 $-\eta_{\mu\nu}$ 部分：
\begin{align*}
 -\eta_{\mu\nu}
 \Tr(\slashed p\gamma^\mu\slashed k\gamma^\nu)
 &=8p\cdot k
 =\frac{4\kappa_\perp^2}{z(1-z)}.
\end{align*}
再收缩轴向规范部分。由 Sudakov 分解，
\begin{equation*}
 P\cdot p=\frac{\kappa_\perp^2}{2z},\qquad
 P\cdot k=\frac{\kappa_\perp^2}{2(1-z)},\qquad
 n\cdot p=zP\cdot n,\qquad
 n\cdot k=(1-z)P\cdot n.
\end{equation*}
因此
\begin{align*}
 &\frac{P_\mu n_\nu+n_\mu P_\nu}{P\cdot n}
 \Tr(\slashed p\gamma^\mu\slashed k\gamma^\nu)\\
 &\quad=\frac{8}{P\cdot n}
 \left[(P\cdot p)(n\cdot k)
 +(P\cdot k)(n\cdot p)
 -(P\cdot n)(p\cdot k)\right]\\
 &\quad=4\kappa_\perp^2
 \left[\frac{1-z}{z}+\frac{z}{1-z}
 -\frac1{z(1-z)}\right]\\
 &\quad=-8\kappa_\perp^2.
\end{align*}
两部分相加并乘母偏振平均 $1/2$ 后，
\begin{align*}
 \mathcal N_{q\bar q}
 &=\frac{T_F}{2}
 \left[
 \frac{4\kappa_\perp^2}{z(1-z)}
 -8\kappa_\perp^2
 \right]\\
 &=2T_F\frac{\kappa_\perp^2}{z(1-z)}
 \left[z^2+(1-z)^2\right].
\end{align*}
近共线母传播子给出 $1/s_{pk}$ 奇点。把完整平方振幅比值严格分解为奇异部分与余项，
\begin{equation*}
 \frac{\overline{|\mathcal M_{g\to q\bar q}|^2}}
 {\overline{|\mathcal M_{g}|^2}}
 \equiv
 \mathcal S_{g\to q\bar q}(z,\kappa_\perp)
 +\mathcal R_{\rm coll}(z,\kappa_\perp),
\end{equation*}
并用近共线极限定义奇异部分
\begin{align*}
 \mathcal S_{g\to q\bar q}
 &\equiv
 g_s^2\frac{\mathcal N_{q\bar q}}{s_{pk}^2}\\
 &=2g_s^2T_F
 \frac{z(1-z)}{\kappa_\perp^2}
 \left[z^2+(1-z)^2\right],
\end{align*}
其中余项满足 $\lim_{\kappa_\perp\to0}\kappa_\perp^2\mathcal R_{\rm coll}=0$。因此分裂核由完整振幅中唯一的 $1/\kappa_\perp^2$ 领先奇点系数确定。
相空间因子也从 Lorentz 不变递推式得到。把两个近共线子粒子先合并成虚母动量 $Q=p+k$，则
\begin{equation*}
 \dd\Phi_{n+1}
 =\dd\Phi_n(\ldots,Q)\,
 \frac{\dd s_{pk}}{2\pi}\,
 \dd\Phi_2(Q;p,k).
\end{equation*}
在 $Q$ 的瞬时静止系中，对两个无质量子粒子有
\begin{equation*}
 \dd\Phi_2
 =\frac1{32\pi^2}\dd\Omega.
\end{equation*}
把方位角积分掉，并用 $z=(1+\cos\theta)/2$，即 $\dd\cos\theta=2\dd z$，得到
\begin{equation*}
 \int\dd\varphi\,\dd\Phi_2
 =\frac1{8\pi}\dd z.
\end{equation*}
因此每个近共线分裂相对于硬 $n$ 体相空间多出
\begin{align*}
 \dd\Phi_{\rm coll}
 &=\frac{\dd s_{pk}}{2\pi}\,
 \frac{\dd z}{8\pi}\\
 &=\frac1{16\pi^2}\dd s_{pk}\,\dd z.
\end{align*}
由
\begin{equation*}
 s_{pk}=\frac{\kappa_\perp^2}{z(1-z)}
\end{equation*}
并在固定 $z$ 的相空间坐标中取微分，
\begin{equation*}
 \dd s_{pk}
 =\frac{\dd\kappa_\perp^2}{z(1-z)},
\end{equation*}
所以
\begin{equation*}
 \dd\Phi_{\rm coll}
 =\frac{1}{16\pi^2}
 \frac{\dd z\,\dd\kappa_\perp^2}{z(1-z)}.
\end{equation*}
于是
\begin{align*}
 \dd P_{g\to q\bar q}
 &=\frac{\overline{|\mathcal M_{g\to q\bar q}|^2}}
 {\overline{|\mathcal M_{g}|^2}}
 \dd\Phi_{\rm coll}\\
 &=\frac{\alpha_s}{2\pi}
 T_F\left[z^2+(1-z)^2\right]
 \frac{\dd\kappa_\perp^2}{\kappa_\perp^2}\,\dd z,
\end{align*}
其中 $g_s^2=4\pi\alpha_s$。因此
\begin{equation*}
 P_{qg}(z)=T_F\left[z^2+(1-z)^2\right].
\end{equation*}

再算 $g(P)\to g(p)g(k)$。选取二维欧几里得横向基 $e_i^\mu$，满足
\begin{equation*}
 e_i\cdot e_j=-\delta_{ij},
 \qquad e_i\cdot P=e_i\cdot n=0,
 \qquad k_\perp^\mu=\kappa_i e_i^\mu.
\end{equation*}
母胶子的物理偏振可取 $\epsilon_0^\mu=e_i^\mu$。两个子胶子的横向基矢必须补上沿 $n^\mu$ 的分量才能同时满足轴向规范条件和各自的质量壳横向条件。逐一解
\begin{equation*}
 n\cdot\epsilon_1=p\cdot\epsilon_1=0,
 \qquad
 n\cdot\epsilon_2=k\cdot\epsilon_2=0
\end{equation*}
得到
\begin{align*}
 \epsilon_1^\mu(a)
 &=e_a^\mu+\frac{\kappa_a}{zP\cdot n}\,n^\mu,\\
 \epsilon_2^\mu(b)
 &=e_b^\mu-\frac{\kappa_b}{(1-z)P\cdot n}\,n^\mu.
\end{align*}
例如
\begin{align*}
 p\cdot\epsilon_1(a)
 &=k_\perp\cdot e_a
 +\frac{\kappa_a}{zP\cdot n}(zP\cdot n)\\
 &=-\kappa_a+\kappa_a=0,
\end{align*}
第二个子胶子同理。

三胶子顶角全部动量取流入。令 $Q=p+k$ 为近在壳母动量，并取
\begin{equation*}
 Q_0=-Q,
 \qquad Q_1=p,
 \qquad Q_2=k,
 \qquad Q_0+Q_1+Q_2=0.
\end{equation*}
去掉共同的 $-g_sf^{abc}$ 后，Lorentz 张量是
\begin{equation*}
 V_{\rho\mu\nu}
 =\eta_{\rho\mu}(Q_0-Q_1)_\nu
 +\eta_{\mu\nu}(Q_1-Q_2)_\rho
 +\eta_{\nu\rho}(Q_2-Q_0)_\mu.
\end{equation*}
定义收缩后的横向分子
\begin{equation*}
 \mathcal V_{iab}
 \equiv
 \epsilon_0^\rho(i)\epsilon_1^\mu(a)\epsilon_2^\nu(b)
 V_{\rho\mu\nu}.
\end{equation*}
三个张量项分别计算。共线领先阶可用 $Q=P+O(\kappa_\perp^2)$。第一项给
\begin{align*}
 (\epsilon_0\cdot\epsilon_1)
 (Q_0-Q_1)\cdot\epsilon_2
 &=-\delta_{ia}
 \left[-(1+z)P-k_\perp\right]\cdot
 \left[e_b-\frac{\kappa_b}{(1-z)P\cdot n}n\right]\\
 &=-\frac{2\delta_{ia}\kappa_b}{1-z}.
\end{align*}
第二项给
\begin{align*}
 (\epsilon_1\cdot\epsilon_2)
 (Q_1-Q_2)\cdot\epsilon_0
 &=(-\delta_{ab})
 \left[(2z-1)P+2k_\perp\right]\cdot e_i\\
 &=2\delta_{ab}\kappa_i.
\end{align*}
第三项给
\begin{align*}
 (\epsilon_2\cdot\epsilon_0)
 (Q_2-Q_0)\cdot\epsilon_1
 &=-\delta_{bi}
 \left[(2-z)P-k_\perp\right]\cdot
 \left[e_a+\frac{\kappa_a}{zP\cdot n}n\right]\\
 &=-\frac{2\delta_{ib}\kappa_a}{z}.
\end{align*}
三项相加，才得到
\begin{equation*}
 \mathcal V_{iab}
 =2\left[
 \delta_{ab}\kappa_i
 -\frac{\delta_{ib}\kappa_a}{z}
 -\frac{\delta_{ia}\kappa_b}{1-z}
 \right]
 +O\!\left(\frac{\kappa_\perp^2}{Q_h}\right).
\end{equation*}

颜色平均为
\begin{align*}
 \frac1{N_c^2-1}\sum_{a,b,c}f^{abc}f^{abc}
 &=\frac{C_A\delta^{aa}}{N_c^2-1}
 =C_A.
\end{align*}
定义
\begin{equation*}
 B_{iab}
 \equiv
 \delta_{ab}\kappa_i
 -\frac{\delta_{ib}\kappa_a}{z}
 -\frac{\delta_{ia}\kappa_b}{1-z},
\end{equation*}
则对两个子偏振求和并对两个母偏振平均需要计算
\begin{equation*}
 \frac12\sum_{i,a,b}\mathcal V_{iab}\mathcal V_{iab}
 =2\sum_{i,a,b}B_{iab}B_{iab}.
\end{equation*}
六类平方与交叉项分别为
\begin{align*}
 \sum_{i,a,b}(\delta_{ab}\kappa_i)^2
 &=2\kappa_\perp^2,\\
 \sum_{i,a,b}\left(\frac{\delta_{ib}\kappa_a}{z}\right)^2
 &=\frac{2\kappa_\perp^2}{z^2},\\
 \sum_{i,a,b}\left(\frac{\delta_{ia}\kappa_b}{1-z}\right)^2
 &=\frac{2\kappa_\perp^2}{(1-z)^2},\\
 -\frac{2}{z}
 \sum_{i,a,b}\delta_{ab}\delta_{ib}\kappa_\ii\kappa_a
 &=-\frac{2\kappa_\perp^2}{z},\\
 -\frac{2}{1-z}
 \sum_{i,a,b}\delta_{ab}\delta_{ia}\kappa_\ii\kappa_b
 &=-\frac{2\kappa_\perp^2}{1-z},\\
 \frac{2}{z(1-z)}
 \sum_{i,a,b}\delta_{ib}\delta_{ia}\kappa_a\kappa_b
 &=\frac{2\kappa_\perp^2}{z(1-z)}.
\end{align*}
因此
\begin{align*}
 \sum_{i,a,b}B_{iab}B_{iab}
 &=2\kappa_\perp^2
 \left[
 1+\frac1{z^2}+\frac1{(1-z)^2}
 -\frac1z-\frac1{1-z}+\frac1{z(1-z)}
 \right]\\
 &=2\kappa_\perp^2
 \frac{(1-z+z^2)^2}{z^2(1-z)^2},
\end{align*}
从而
\begin{equation*}
 \frac12\sum_{i,a,b}\mathcal V_{iab}\mathcal V_{iab}
 =4\kappa_\perp^2
 \frac{(1-z+z^2)^2}{z^2(1-z)^2}.
\end{equation*}
除以母传播子虚度平方 $s_{pk}^2=\kappa_\perp^4/[z^2(1-z)^2]$，再乘同一个共线相空间，得到
\begin{align*}
 \dd P_{g\to gg}
 &=\frac{\alpha_s}{2\pi}
 \left\{2C_A
 \frac{(1-z+z^2)^2}{z(1-z)}\right\}
 \frac{\dd\kappa_\perp^2}{\kappa_\perp^2}\,\dd z.
\end{align*}
最后使用恒等式
\begin{align*}
 \frac{(1-z+z^2)^2}{z(1-z)}
 &=\frac{z}{1-z}+\frac{1-z}{z}+z(1-z)
\end{align*}
得到实辐射核
\begin{equation*}
 P_{gg}^{\rm real}(z)
 =2C_A\left[
 \frac{z}{1-z}+\frac{1-z}{z}+z(1-z)
 \right].
\end{equation*}

最后，$q\to qg$ 的实辐射已经在正文\eqnrefs{eq:q-to-qg-splitting-probability-dp10} 中推得。若现在把被观测子粒子改成胶子，令其携带 $z_g=1-z_q$，则
\begin{align*}
 P_{gq}(z_g)
 &=P_{qq}^{\rm real}(1-z_g)\\
 &=C_F\frac{1+(1-z_g)^2}{z_g}.
\end{align*}
因此三个实辐射核全部从同一共线母式得到。$P_{gg}$ 的 $z\to1$ 端点必须与虚修正共同解释为加号分布，而 $\delta(1-z)$ 系数由夸克数与总动量守恒固定；合并这些端点项后得到正文\eqnrefs{eq:qcd-lo-splitting-kernels-rest-dp83}。
\end{derivation}

\begin{derivation}[LO DGLAP 守恒和规则]
先从 DGLAP 方程的矩结构说明为什么需要端点项。若 $f_i(x,\mu)$ 是部分子分布，则
\begin{equation*}
 \frac{\partial f_i(x,\mu)}{\partial\ln\mu^2}
 =\frac{\alpha_s}{2\pi}
 \sum_j\int_x^1\frac{\dd z}{z}
 P_{ij}(z)f_j\!\left(\frac{x}{z},\mu\right).
\end{equation*}
对非单态夸克组合积分 $\int_0^1\dd x$，交换 $x,z$ 积分并令 $y=x/z$：
\begin{align*}
 \frac{\dd}{\dd\ln\mu^2}\int_0^1\dd x\,q_{\rm NS}(x)
 &=\frac{\alpha_s}{2\pi}
 \int_0^1\dd z\,P_{qq}(z)
 \int_0^1\dd y\,q_{\rm NS}(y).
\end{align*}
非单态夸克数守恒因此要求
\begin{equation*}
 \int_0^1\dd z\,P_{qq}(z)=0,
\end{equation*}
这正是正文中 $P_{qq}$ 的 $3\delta(1-z)/2$ 项来源。

总纵向动量分数为
\begin{equation*}
 M(\mu)\equiv\sum_i\int_0^1\dd x\,x f_i(x,\mu).
\end{equation*}
同样对 DGLAP 方程乘 $x$ 并积分，令 $x=zy$，得到
\begin{align*}
 \frac{\dd M}{\dd\ln\mu^2}
 &=\frac{\alpha_s}{2\pi}
 \sum_{i,j}\int_0^1\dd z\,zP_{ij}(z)
 \int_0^1\dd y\,y f_j(y,\mu).
\end{align*}
由于不同母部分子 $j$ 的分布彼此独立，要使任意强子态都满足 $\dd M/\dd\ln\mu^2=0$，必须逐个母部分子满足
\begin{equation*}
 \sum_i\int_0^1\dd z\,zP_{ij}(z)=0.
\end{equation*}

对母夸克 $j=q$，先算 $P_{qq}$ 的一阶动量矩。加号分布定义给
\begin{align*}
 \int_0^1\dd z\,
 z\frac{1+z^2}{(1-z)_+}
 &=\int_0^1\dd z\,
 \frac{z(1+z^2)-2}{1-z}\\
 &=\int_0^1\dd z\,
 \frac{z+z^3-2}{1-z}.
\end{align*}
多项式除法
\begin{equation*}
 z+z^3-2=-(1-z)(z^2+z+2)
\end{equation*}
于是
\begin{align*}
 \int_0^1\dd z\,
 z\frac{1+z^2}{(1-z)_+}
 &=-\int_0^1\dd z\,(z^2+z+2)\\
 &=-\frac13-\frac12-2
 =-\frac{17}{6}.
\end{align*}
加上虚修正得到
\begin{align*}
 \int_0^1\dd z\,zP_{qq}(z)
 &=C_F\left(-\frac{17}{6}+\frac32\right)
 =-\frac43C_F.
\end{align*}
另一方面
\begin{align*}
 \int_0^1\dd z\,zP_{gq}(z)
 &=C_F\int_0^1\dd z\,
 \left[1+(1-z)^2\right]\\
 &=C_F\int_0^1\dd z\,(2-2z+z^2)\\
 &=C_F\left(2-1+\frac13\right)
 =\frac43C_F.
\end{align*}
两者相加为零，母夸克的动量和规则成立。

对母胶子 $j=g$，一个给定夸克味贡献
\begin{align*}
 \int_0^1\dd z\,zP_{qg}(z)
 &=T_F\int_0^1\dd z\,
 z\left[z^2+(1-z)^2\right]\\
 &=T_F\int_0^1\dd z\,
 (2z^3-2z^2+z)\\
 &=T_F\left(\frac12-\frac23+\frac12\right)
 =\frac{T_F}{3}.
\end{align*}
夸克和反夸克都要计入，因此 $n_f$ 个活跃味总共给 $2n_fT_F/3$。

把胶子核的虚修正暂写成 $A_g\delta(1-z)$。实辐射的三个动量矩逐项为
\begin{align*}
 I_1
 &\equiv\int_0^1\dd z\,
 z\frac{z}{(1-z)_+}
 =\int_0^1\dd z\,
 \frac{z^2-1}{1-z}\\
 &=-\int_0^1\dd z\,(1+z)
 =-\frac32,\\[2mm]
 I_2
 &\equiv\int_0^1\dd z\,z\frac{1-z}{z}
 =\int_0^1\dd z\,(1-z)
 =\frac12,\\[2mm]
 I_3
 &\equiv\int_0^1\dd z\,z^2(1-z)
 =\frac13-\frac14
 =\frac1{12}.
\end{align*}
因此
\begin{align*}
 \int_0^1\dd z\,zP_{gg}(z)
 &=2C_A(I_1+I_2+I_3)+A_g\\
 &=-\frac{11}{6}C_A+A_g.
\end{align*}
母胶子动量守恒要求
\begin{equation*}
 -\frac{11}{6}C_A+A_g+\frac23T_Fn_f=0,
\end{equation*}
所以
\begin{align*}
 A_g
 &=\frac{11}{6}C_A-\frac23T_Fn_f\\
 &=\frac12\left(\frac{11}{3}C_A-\frac43T_Fn_f\right)
 =\frac{\beta_0}{2}.
\end{align*}
\end{derivation}
